% Part: lambda-calculus
% Chapter: introduction
% Section: fixed-point-combinator

\documentclass[../../../include/open-logic-section]{subfiles}

\begin{document}

\olfileid{lam}{int}{fix}
\olsection{స్థిరబిందు సంయోజకాలు}

మీ దగ్గర లాంబ్డా పదం~$g$ ఉందనీ, మరొక పదం~$k$ కావాలనీ, ఆ $k$ అనేది
$gk$కు $\beta$-తుల్యం కావాలనీ అనుకుందాం. మన సంకేత సంప్రదాయాలతో
\[
\fn{diag}(x) = xx
\]
మరియు
\[
l(x) = g(\fn{diag}(x))
\]
అనే పదాలను నిర్వచించండి; మరో మాటలో, $l$ అనేది
$\lambd[x][g(xx)]$ అనే పదం. $k$ను $ll$ అనే పదంగా ఉంచండి. అప్పుడు
\begin{align*}
k & = (\lambd[x][g(xx)])(\lambd[x][g(xx)]) \\
& \red  g((\lambd[x][g(xx)])(\lambd[x][g(xx)])) \\
& = gk.
\end{align*}
\[
Y = \lambd[g][((\lambd[x][g(xx)])(\lambd[x][g(xx)]))]
\]
అని తీసుకుంటే $Yg$, $g(Yg)$ రెండూ ఒకే పదానికి తగ్గుతాయి; కనుక
$Yg \equiv_\beta g(Yg)$. దీన్ని ``కరీ సంయోజకం'' అంటారు. బదులుగా
\[
Y = (\lambd[xg][g(xxg)])(\lambd[xg][g(xxg)])
\]
అని తీసుకుంటే, నిజానికి $Yg$ అనేది $g(Yg)$కు తగ్గుతుంది; ఇది మరింత
బలమైన ప్రకటన. $Y$ యొక్క ఈ రెండవ రూపాన్ని ``ట్యూరింగ్ సంయోజకం'' అంటారు.

\end{document}
